% -------------------------------------------------------------------------
% WBS ID: RES-STR-DMG
% Deliverable: Damage, Failure and Performance-Loss Models
% Status: In progress
% Evidence status: Damage model deferred; load requirements identified
% Review status: Not reviewed
% Last updated: 2026-07-19
% Dependencies: RES-STR-LOD, RES-STR-RSP, RES-VER-MET
% Downstream interfaces: Future structural damage extension
% -------------------------------------------------------------------------

\section{RES-STR-DMG --- Damage, Failure and Performance-Loss Models}
\label{sec:res-str-dmg}

\subsection{Purpose and Scope}
\label{sec:res-str-dmg-purpose}

This Deliverable records damage and failure modelling as a stretch
extension. The active project baseline is hydrodynamic comparison of
fixed rigid barriers; it shall not infer material damage, scour,
fracture or post-failure debris motion from fluid loads alone.

\subsection{Research Questions}
\label{sec:res-str-dmg-research-questions}

% TODO: State the questions that must be answered before the Deliverable
% can be accepted.

\subsection{Terminology and Definitions}
\label{sec:res-str-dmg-definitions}

% TODO: Define the technical terminology, symbols, quantities and
% classifications used in this Deliverable.

\subsection{Literature Evidence}
\label{sec:res-str-dmg-literature-evidence}

% TODO: Record evidence from peer-reviewed papers, authoritative datasets,
% standards, technical reports and primary documentation.
%
% Do not create a detached literature-summary list. Explain how each source
% supports, contradicts or limits a project decision.

\subsection{Governing Relations and Quantitative Definitions}
\label{sec:res-str-dmg-governing-relations}

% TODO: Record relevant equations, dimensionless groups, empirical models,
% extraction rules or quantitative definitions.
%
% Use align environments for displayed equations.
% Every displayed equation must have a unique \label{eq:...}.
% Do not place punctuation after displayed equations.

\subsection{Material or Structural System}
\label{sec:res-str-dmg-material-structural-system}

% TODO: Complete this RES-STR Domain-specific subsection.

\subsection{Load Cases}
\label{sec:res-str-dmg-load-cases}

% TODO: Complete this RES-STR Domain-specific subsection.

\subsection{Constitutive Behaviour}
\label{sec:res-str-dmg-constitutive-behaviour}

% TODO: Complete this RES-STR Domain-specific subsection.

\subsection{Response and Stability Measures}
\label{sec:res-str-dmg-response-stability-measures}

% TODO: Complete this RES-STR Domain-specific subsection.

\subsection{Damage and Failure Modes}
\label{sec:res-str-dmg-damage-failure-modes}

% TODO: Complete this RES-STR Domain-specific subsection.

\subsection{Initial Failure-Modelling Scope}
\label{subsec:res_str_dmg_initial_failure_scope}

The Updated-Lagrangian formulation permits large deformation but does not independently provide plasticity, damage, fracture or topology change. Earlier structural notes identified this as a possible future route; it is not the active baseline.

Full post-fracture continuation would additionally require crack or damage evolution, contact between separated surfaces, topology change and a fluid treatment for detached structural fragments. These capabilities are outside the active fixed-rigid hydrodynamic baseline.

\textbf{Selected scope.} Defer large nonlinear deformation, material
damage, fragmentation, debris--fluid interaction and scour until the
fixed-rigid hydrodynamic load histories have been validated and
accepted. Source role: \textcite{baragamage2024fullycoupled}
demonstrates that tsunami loading can be coupled to structural motion,
but also supports treating such coupling as a higher-fidelity
extension. Limitation: no project damage model has been selected or
executed.


\subsection{Fluid--Structure Coupling Requirements}
\label{sec:res-str-dmg-fluid-structure-coupling-requirements}

% TODO: Complete this RES-STR Domain-specific subsection.

\subsection{Candidate Approaches}
\label{sec:res-str-dmg-candidate-approaches}

% TODO: Define the credible candidate approaches considered.

\subsection{Critical Comparison}
\label{sec:res-str-dmg-critical-comparison}

% TODO: Compare candidates using physically and project-relevant criteria.
% Do not select an approach solely on popularity or implementation ease.

\subsection{Assumptions and Applicability}
\label{sec:res-str-dmg-assumptions}

% TODO: State assumptions and the conditions under which they are valid.

\subsection{Limitations and Failure Conditions}
\label{sec:res-str-dmg-limitations}

% TODO: State where the evidence, model or selected approach ceases to be
% reliable or applicable.

\subsection{Project-Specific Conclusions}
\label{sec:res-str-dmg-conclusions}

The current Studentship output is a load dataset suitable for later
damage studies, not a damage prediction. Damage claims require a
separate constitutive model, validation evidence and acceptance gate.

\subsection{Selected Approach and Rationale}
\label{sec:res-str-dmg-selected-approach}

The selected approach is formal deferral. Damage metrics may be listed
as future consumers of pressure, shear, force, impulse and moment
histories, but shall not be reported as baseline simulation outputs.

\subsection{Unresolved Decisions}
\label{sec:res-str-dmg-unresolved-decisions}

Open decisions are candidate material families, damage variables,
failure criteria, scour coupling and validation data for damage or
post-failure behaviour.

\subsection{Requirements and Handoffs}
\label{sec:res-str-dmg-requirements}

Software shall not implement damage or deformable-structure logic from
this deliverable. It shall preserve hydrodynamic histories and
geometry/patch identifiers so a later structural extension can consume
them.

\subsection{Deliverable Completion Record}
\label{sec:res-str-dmg-completion-record}

% TODO: Record whether the Deliverable is:
%
% - Not started
% - In progress
% - Draft complete
% - Reviewed
% - Accepted
%
% Acceptance requires:
%
% - the research questions to be answered;
% - claims to be supported by appropriate evidence;
% - assumptions and limitations to be explicit;
% - the project-specific conclusion to be stated;
% - downstream requirements to be traceable;
% - unresolved decisions to be closed or formally deferred.
